% sections/introduction.tex

\subsection{The Problem of Lineage Commitment}

Among the most consequential decisions a cell makes is the irreversible commitment to a particular developmental lineage. In the blood-forming system, a common myeloid progenitor must resolve between the erythroid fate---producing oxygen-carrying red blood cells---and the myeloid fate, giving rise to the granulocytes and monocytes of the innate immune system. This decision is regulated by cytokines: erythropoietin (EPO) promotes erythroid commitment, and granulocyte colony-stimulating factor (GCSF) promotes myeloid commitment. The molecular arbiters of the decision are two master transcription factors, GATA1 and PU.1, which antagonise each other through a double-negative feedback loop. When GATA1 dominates, erythroid genes are expressed and PU.1 activity is suppressed; when PU.1 dominates, the myeloid programme proceeds and GATA1 is degraded. The architecture of this toggle switch has been known since the mid-1990s, yet a fundamental question about how it operates has resisted resolution for more than thirty years: does the cytokine environment instruct fate by deterministically driving the switch, or does the cell make a stochastic choice that cytokines merely bias in probability?

\subsection{Three Decades of Unresolved Debate}

The question was first sharpened by Metcalf~\cite{metcalf1988}, who observed from bulk colony assays that specific cytokines reliably produced specific blood cell lineages, suggesting instruction. Cantor and Orkin~\cite{cantor2002} later described the transcriptional architecture of erythropoiesis in mechanistic terms that reinforced the instructive view: EPO activates GATA1 through receptor-coupled signalling, GATA1 suppresses PU.1 through direct interaction, and the outcome is therefore predictable from the input. Yet when single-cell and clonal experiments became possible, the picture changed substantially. Enver \etal~\cite{enver1998} noted that progenitors showed fate heterogeneity that could not be explained by instructive models, asking whether stem cells play dice. Hoppe \etal~\cite{hoppe2016} resolved individual cell fate by live imaging of GATA1-GFP and PU.1-mCherry reporters in clonal populations under identical cytokine exposure and found that cell fate was determined before any EPO-dependent protein ratio divergence at the population level---a result incompatible with pure instruction. Olsson \etal~\cite{olsson2016} used single-cell RNA sequencing to trace progenitor trajectories and observed that mixed-lineage states preceded binary commitment, consistent with a system that commits stochastically from a primed intermediate. Meanwhile, Chang \etal~\cite{chang2008} demonstrated in a different progenitor context that transcriptome-wide noise quantitatively predicted lineage choice, directly implicating molecular stochasticity in fate decisions.

Both bodies of evidence are internally consistent. The apparent contradiction between them has driven thirty years of parallel experimental traditions---bulk assays producing instructive-looking data and single-cell assays producing stochastic-looking data from the same biological system. A mechanistic explanation for how both observations can simultaneously be true has not been provided.

\subsection{Computational Modelling of the Toggle Switch}

The GATA1/PU.1 system has attracted substantial modelling attention. Huang \etal~\cite{huang2007} constructed an ordinary differential equation model of the double-negative feedback and demonstrated that the circuit produces bistability: two stable attractors corresponding to erythroid and myeloid committed states, separated by an unstable equilibrium. Their model captured the qualitative phenomenology of commitment but, being deterministic, could not address fate variability at single-cell copy numbers. Chickarmane \etal~\cite{chickarmane2009} extended this work by introducing cooperativity in the cross-repression terms and analysing basin geometry, identifying primed-state residence time as a key predictor of commitment delay. These models established the bistable framework within which the current debate operates, but they address population-level attractors, not the single-cell probability of entering each attractor.

What has been missing is a model that explicitly treats molecular copy number as a variable, spans the range from single-cell to population scale, and decomposes the dynamics into mechanistic layers---reception, decision, and execution---so that the role of noise can be localised within the circuit. Stochastic modelling of gene regulatory networks is well-established through Gillespie-based simulation~\cite{gillespie1977}, but its application to the full hematopoietic commitment circuit, including receptor kinetics, energy metabolism, and the phosphorylation cascade that locks in commitment, has not been carried out in a unified framework.

\subsection{Approach and Contributions}

In this work we construct a stochastic Petri net model of the GATA1/PU.1 commitment circuit implemented in Signal Hierarchical Petri Nets~\cite{shypn2026}. The model integrates cytokine receptor binding and recycling, mRNA transcription and nuclear export, protein translation, nuclear import, phosphorylation, and ubiquitin-mediated degradation, together with the energy metabolism (ATP, GTP) that powers these transitions. Consistent with the established biology that uncommitted progenitors carry a slight GATA1 bias before cytokine exposure~\cite{cantor2002}, we deliberately encode a 33\% kinetic advantage for GATA1 transcription (basal rate 0.08~mM/s versus 0.06~mM/s for PU.1); this asymmetry is a structural modelling assumption, not a calibration artefact. By scaling initial molecular concentrations to 0.01$\times$ standard values, we enter the regime of physiological single-cell copy numbers, where stochastic fluctuations dominate.

We perform N=50 independent stochastic simulations at each of eight EPO concentrations (Phase~B) and across ten GCSF conditions at fixed EPO (Phase~D and Phase~D-ext), analysing both ensemble statistics and mean trajectories. Our central findings are these. Phase~B commitment probability is 60--70\% and statistically insensitive to EPO dose alone, demonstrating that stochastic noise---not cytokine concentration---is the proximate determinant of individual cell fate. Phase~D and D-ext reveal the true bifurcation structure: P(erythroid) is $\geq$\,0.92 for GCSF/EPO\,$\leq$\,2.2 and switches abruptly to zero at GCSF/EPO\,=\,2.5, locating the commitment threshold at a GCSF/EPO ratio of $\approx$\,2.3 rather than at an EPO dose. Analysis of the mean trajectory reveals a three-layer cascade architecture: a reception layer whose output is structurally fixed independent of EPO dose, a decision layer in which the fate-determining fluctuations arise, and an execution layer that is rapid and EPO-dose-insensitive once triggered. The $t\approx146$~s GATA1/PU.1 priming crossover establishes cytokine-dependent GATA1 dominance---not an irreversible molecular record; a Phase~X cytokine-withdrawal experiment demonstrates complete fate reversal to the myeloid attractor within $\approx$\,1500~s when EPO is removed. The EPO* threshold visible in population assays is a statistical artefact of averaging that disappears at single-cell resolution.

These findings reconcile the instructive and stochastic views mechanistically and generate five experimentally testable predictions that provide rational guidance for future single-cell wet-lab experiments, including a precisely bounded EPO pulse-duration window and a direct test for the population-averaging artefact.
